% §7 Discussion
We are explicit about what this paper does and does not claim.

\paragraph{Scope of the evidence.}
The formal results are the contribution; the orchestrator PoC verifies the \emph{semantics}
the results assume, and the smoke anchor instantiates the quantities on a simulated delivery
clock.
There is no real ASR, TTS, or language model in the loop, and the frame and tool latencies are
synthetic.
The anchor is therefore a feasibility instance with arbitrary units, not a benchmark.
The per-architecture cessation values a production system would observe also depend on its
specific synthesis queueing and its backend tool latency, which real deployment telemetry
confirms can reach $\sim$1~s per round trip.

\paragraph{No uniqueness claim.}
The event-driven configuration is presented as a \emph{sufficient}, idiomatic way to own both
floors (Definitions~\ref{def:clocks} and~\ref{def:interruptibility}).
Serial topologies can host async nodes, background dispatch, and checkpoint reentry and thereby
emulate the same capability; the paper's claim is about idiomatic expressiveness, not
in-principle precedence.
Similarly, Proposition~\ref{prop:floor_dependency} is a dependency statement under a fixed in-flight
policy, not a claim that the floors are independent in general.

\paragraph{Policy vs physics.}
Theorem~\ref{thm:cessation_floor} separates \emph{can stop} from \emph{should stop}.
Ceasing fast is not always the right user experience, and an engagement-aware policy may
legitimately decide to keep the floor \citep{bohus2011interruptions}.
The floor is a capability bound, and the paper makes no claim about optimising raw cession speed
in every turn.

\paragraph{Relation to product vocabulary.}
The onset side maps one-to-one onto the server-side voice-to-voice metric used in deployment --
the interval from detected end of user speech to first audio leaving the synthesis server.
This paper adds the symmetrical interval on the other side of an interruption, which
production telemetry today does not emit at all.
Of the two, cessation is the one most frequently felt on tool-heavy turns and the one with no
standard definition, reinforcing the motivation of Section~\ref{sec:introduction}.

\paragraph{Citation status.}
Related-work references were verified against the arXiv API and Crossref on 2026-09-01,
including a 2025--26 sweep that surfaced the full-duplex and interruption benchmarks cited
in Section~\ref{sec:background}.
The companion onset-floor paper is available as an alphaXiv preprint
(alphaXiv:2608.latency-floor-cascade-voice-agents); the formal results of
Section~\ref{sec:analysis} do not depend on its indexing status.

\paragraph{Relation to the companion paper.}
Where the companion paper \citep{latencyfloor2025} bounded response onset from outside the stage, this paper bounds
response cessation from inside it, and reuses the same reducible/irreducible intuition and the
same $200$~ms perceptual anchor.
Together the two papers give a two-sided picture of realtime responsiveness for cascade
agents.
